\section{Token Efficiency: el optimizador Muon y la estabilidad del entrenamiento}

\subsection{Por qué la token efficiency es tan importante}

Yang Zhilin (杨植麟) parte de la clásica Kaplan scaling law para enfatizar que la token efficiency no solo concierne a la eficiencia, sino también al \textbf{límite superior de la inteligencia}.

\begin{importantbox}{Token Efficiency = incremento del límite superior de la inteligencia}
Supongamos que se cuenta con 50 billones de tokens de alta calidad; si un nuevo optimizador logra duplicar la token efficiency, eso equivale a obtener «de la nada» el efecto de 100 billones de tokens. En una época en la que el volumen de datos de alta calidad se acerca a su límite (data wall), incrementar la token efficiency significa directamente elevar el techo de la inteligencia.
\end{importantbox}

\subsection{El optimizador Muon}

Muon es un \textbf{optimizador de segundo orden} cuya idea central es transformar cada actualización de gradiente de modo que las distintas componentes de los parámetros queden ortogonales entre sí. En comparación con el tradicional AdamW, Muon mejora de manera notable el desempeño con la misma cantidad de parámetros y de tokens de entrenamiento.

El equipo de Kimi fue el primero en demostrar que el optimizador Muon puede escalar al entrenamiento de LLM a gran escala; las técnicas centrales incluyen:
\begin{itemize}
    \item \textbf{Estrategia de decay}: fundamental para escalar a modelos más grandes
    \item \textbf{Consistencia de RMS}: se introduce un coeficiente ajustable para asegurar que la actualización de RMS de Muon sea comparable con la de Adam
    \item \textbf{Implementación distribuida}: se fragmenta el estado del optimizador entre el data parallel group, lo que logra una implementación distribuida eficiente de Muon
\end{itemize}

\subsection{QK-Clip: cómo resolver la inestabilidad del entrenamiento}

Al escalar Muon a un modelo de 1 billón de parámetros, el equipo se topó con un problema de inestabilidad en el entrenamiento: los max logits se disparaban rápidamente por encima de 1000 (el valor normal es de alrededor de 50 a 100), lo que provocaba que el entrenamiento divergiera.

\begin{importantbox}{La técnica QK-Clip}
Para cada attention head, se calcula el max logit durante el forward pass y luego se calcula un factor divisor que se aplica a las proyecciones de query y de key, con lo que se restringe el valor máximo dentro de un rango dado. Los experimentos muestran que:
\begin{itemize}
    \item Las curvas de entrenamiento antes y después de aplicar QK-Clip coinciden de manera estricta, es decir, no afecta la convergencia del entrenamiento
    \item Una vez que el max logit alcanza el umbral (por ejemplo, 100), queda restringido de manera efectiva y después disminuye de forma natural
\end{itemize}
Esto permitió que el equipo de Kimi completara con éxito el primer entrenamiento a gran escala con Muon en la historia (1 billón de parámetros).
\end{importantbox}

\begin{warningbox}{La estabilidad del entrenamiento no se puede ignorar}
El problema de logit explosion, que rara vez aparece en el entrenamiento tradicional con Adam, puede amplificarse al usar un optimizador de segundo orden. Cambiar el optimizador sin más, sin considerar mecanismos de estabilidad, puede provocar que el entrenamiento fracase por completo.
\end{warningbox}

\subsection{Resumen de la sección}

La token efficiency no es solo una optimización de ingeniería, sino también una vía clave para elevar el límite superior de la inteligencia. El optimizador Muon logra duplicar la token efficiency al ortogonalizar las actualizaciones de gradiente; la técnica QK-Clip resuelve el problema de estabilidad en el entrenamiento a gran escala y permite que Muon escale con éxito hasta el billón de parámetros.

